As introduced in Section~1, AdsMind comprises five modules: \textbf{Planner}, \textbf{Validator}, \textbf{Executor}, \textbf{Analyzer}, and \textbf{Summarizer}.
We now detail their implementation.

\subsection{Closed-Loop Agent Architecture}

\subsubsection{Planner and Validator: Adsorption Hypothesis Generation}

The user provides a request prompt, adsorbate SMILES string, and surface slab file.
At each iteration, the \textbf{Planner} proposes a structured adsorption hypothesis in JSON format, including the canonical site type (\textit{ontop}, \textit{bridge}, or \textit{hollow}), intended surface symbols, and adsorbate binding indices~\cite{norskov2004density,shriver2006inorganic}.

In the closed-loop architecture, updated context from other modules is entered into the \textbf{Planner} prompt as physical feedback.
The prompt is assembled from conditional sections controlled by three runtime switches:
\textit{Slip feedback} frames prior slips as evidence that the planned site is unstable;
\textit{Forbid} converts slipped or failed site intents into explicit prohibitions;
\textit{Termination} surfaces convergence-to-known-best signals and allows early stopping.
The three switches are independent; disabling all three yields the 1-Shot baseline (loop continues until attempt budget is exhausted, default: 5).

Convergence and stopping criteria are managed by the \textbf{Summarizer} (see below).

After the hypothesis is proposed, the \textbf{Validator} enforces chemistry- and schema-level constraints before any expensive physics simulation:
(1)~requested adsorbate binding indices are valid and correspond to non-hydrogen heavy atoms (unless the adsorbate is a single atom);
(2)~cardinality of binding indices matches the site type (e.g., \textit{ontop} requires exactly one index);
(3)~the hypothesis has not been previously evaluated, preventing redundant calculations.
Failed hypotheses elicit an immediate penalty feedback, directing the LLM to refine its proposal without advancing to physical simulation.

\subsubsection{Executor and Analyzer: Physical Simulation and Scientific Analysis}

The \textbf{Executor} instantiates the validated plan on the slab, performs physical relaxation, and returns structured post-relaxation analysis.
The \textbf{Analyzer} records whether the final structure remained bound, whether dissociation or rearrangement occurred, and whether the realized coordination environment differs from the \textbf{Planner}'s intent.
\textit{Chemical Slip} is detected when the relaxed structure exhibits a different surface-symbol fingerprint (sorted element symbols within the first coordination shell of the anchor atom) than the Planner's intent; canonical site-family mismatch is retained separately as a geometric slip signal.

The proposal is pre-processed through a surrogate-SMILES geometry layer that maps binding atoms to proxy markers and constructs initial placements for \textit{ontop}, \textit{bridge}, and \textit{hollow} adsorption.
It samples a capped set of candidate surface sites and conformers (default: eight filtered coordinates, four conformers per site; one top-ranked candidate proceeds to full relaxation).

A \textbf{two-step evaluation procedure} is employed: candidates are first screened with the MACE-MP calculator to discard pathological configurations and rank by screened energy (Langevin warmup disabled in primary benchmark: \texttt{md\_steps=0}); the best surviving candidate (\texttt{relax\_top\_n=1}) is then fully relaxed and analyzed.

The adsorption energy is:
\begin{equation}
    E_{\text{ads}} = E_{\text{total}} - E_{\text{surface}} - E_{\text{adsorbate}}
\end{equation}
where $E_{\text{total}}$ is the relaxed adsorbate-slab energy, $E_{\text{surface}}$ is the bare surface single-point energy, and $E_{\text{adsorbate}}$ is the isolated adsorbate energy.

Relaxation uses the BFGS algorithm~\cite{nocedal2006numerical} with maximum force below $f_{\text{max}}$:
\begin{equation}
    \max_{i} \| \mathbf{F}_i \|_{\infty} \le f_{\text{max}}
\end{equation}
The bottom third of the slab is fixed; upper layers and adsorbate are fully relaxed.

The \textbf{Analyzer} parses the relaxed trajectory to extract the local coordination environment.
Connectivity is established via an adjacency matrix $A$: atoms $i$ and $j$ are bonded if $d_{ij} \le \gamma(R_i + R_j)$, with $\gamma=1.35$ for intramolecular dissociation checks and $\gamma=1.30$ (relaxed to $1.45$ for $E_{\text{ads}} < -0.5$~eV) for adsorbate--surface coordination.
Dissociation is flagged if the adsorbate subgraph has $>1$ connected component; the coordination environment classifies the site as \textit{ontop}, \textit{bridge}, or \textit{hollow}.
Slip is logged if the realized site or surface-symbol fingerprint differs from the Planner's intent.
The \textbf{Analyzer} updates running memory with adsorption energy, bond integrity, and slip status, enabling dynamic adjustment of the exploration strategy.

\subsubsection{Summarizer: Convergence Monitoring and Early Stopping}

The \textbf{Summarizer} consolidates the running memory maintained by the Analyzer across iterations.
Convergence is declared when the same site family is proposed in three consecutive iterations without finding a lower-energy configuration.
Upon convergence, the Summarizer emits a STOP signal to terminate the search early, preventing unnecessary computation when further exploration is unlikely to improve the result.

\subsection{Physical Backend: MACE-MP-0 Calculator}
\label{sec:mace_config}

The physical engine uses the MACE equivariant message-passing architecture with the MACE-MP foundation model~\cite{batatia2022mace,batatia2023foundation}.
Two configurations are employed:
\textbf{Primary benchmark}: MACE-MP-0 small, CPU, float32, dispersion disabled (conservative setting held fixed across all ablations).
\textbf{Sensitivity checks}: MACE-MP-0 large, CUDA, float64, dispersion enabled (probes force-field sensitivity separately).
All primary benchmark and matched-baseline experiments use the same MACE-MP-0 small backend so that observed differences reflect search strategy, not energy-function choice. Key parameters are in Table~\ref{tab:mace_config}.

\subsection{Evaluation Protocol and Experimental Setup}
\label{sec:eval_setup}

We evaluate five variants in a systematic ablation:
\begin{enumerate}
  \item \textbf{Full}: all mechanisms enabled
  \item \textbf{w/o Slip}: removes slip feedback and FORBID constraints
  \item \textbf{w/o Forbid}: retains slip detection but suppresses prohibition
  \item \textbf{w/o Term}: disables early stopping
  \item \textbf{1-Shot}: single attempt, all mechanisms disabled
\end{enumerate}

All benchmark runs use frozen experiment configurations. Four reasoning backends are evaluated at zero temperature: Gemini~2.5~Pro, Grok-4, GPT-5.4, and Claude~Sonnet~4.6 (see SI for model identifiers).
The physical backend is locked to MACE-MP-0 small (\S~\ref{sec:mace_config}).

\subsection{CMU20 Experimental Setup}

The CMU20 benchmark uses the completed 20-case manifest, spanning H, OH, NNH, CH$_2$CH$_2$OH, OCHCH$_3$, and ONN(CH$_3$)$_2$ across intermetallic and monometallic surfaces.
Full runs use max five attempts; 1-shot restricts to a single attempt.
The five-variant ablation matrix is executed on all four LLM backends, yielding 400 total attempts.
Natural failures (dissociation/reaction) are retained in the success-rate denominator; only external service/calculator errors are excluded.

\subsection{OCD-GMAE62 Benchmark and Two-Tier Evaluation Protocol}

For validation beyond CMU20, we assemble \textbf{OCD-GMAE62}: a 62-case dataset spanning diverse surfaces (intermetallic, monometallic, compound) and adsorbates (H, OH, NNH, CH$_2$CH$_2$OH, OCHCH$_3$, ONN(CH$_3$)$_2$, and others).
To systematically evaluate performance and stability, we employ a \textbf{two-tier protocol}:
Tier~1 is a full-matrix evaluation across all 62 cases and five variants (Full, w/o Slip, w/o Forbid, w/o Term, 1-Shot), providing baseline reliability, 1-Shot penalty quantification, and cross-backend agreement.
Tier~2 is a stability audit on 12 subsampled cases with $N=5$ independent repeated runs each, quantifying run-to-run variance and diagnosing stochasticity.
All experiments use four LLM backends under fixed MACE-MP-0 small CPU protocol; external service failures are excluded, natural failures are retained.

The 62-case OCD-GMAE62 dataset includes 12 cases that were part of a preliminary evaluation conducted prior to the full dataset construction.
These 12 cases were selected for the stability audit because they (1) span diverse surface compositions (intermetallic and monometallic),
(2) include diverse adsorbates (H, OH, NNH, OCHCH$_3$, etc.), and (3) have existing single-run records from the preliminary evaluation,
enabling a direct comparison of run-to-run variability.
The remaining 50 cases---together with 10 of the 12 that have single-run records---form the Tier~1 full-matrix cohort (one run per case per backend).
The 12 Tier~2 cases are a subset of the 62-case dataset, not an independent collection; the preliminary evaluation records are used only for reproducibility audit, not for performance evaluation.
